\documentclass[a4paper,12pt]{article}
\usepackage[utf8]{inputenc}
\usepackage[T1]{fontenc}
\usepackage[english]{babel}
\usepackage{geometry}
\geometry{left=2.5cm,right=2.5cm,top=2.5cm,bottom=2.5cm}

\usepackage{lmodern}
\usepackage{xcolor}
\definecolor{titleblue}{RGB}{0,51,140}

\usepackage{hyperref}

\pagestyle{empty}

\begin{document}

\begin{flushleft}
    \textbf{Ismail AISSAOUI} \\
    State Engineer in Artificial Intelligence \\
    ismail.aissaoui.pro@gmail.com \\
    +213 660 70 77 96 \\
    \href{https://ismail-aissaoui.vercel.app}{ismail-aissaoui.vercel.app}
\href{https://linkedin.com/in/aissaoui-ismail-6a77a92a7}{linkedin.com/in/aissaoui-ismail-6a77a92a7}
\end{flushleft}

\begin{flushright}
    To the PhD Selection Committee \\
    \textbf{EDT Program - Catalyst Project} \\
    Attn: Prof. Fabrice Gamboa, Dr. Mathieu Vallee, Dr. Clément Gauchy \\
    \medskip
    April 30, 2026
\end{flushright}

\bigskip

\noindent \textbf{Subject: Application for PhD: Interpretable and Robust Machine Learning for Physics-Based Numerical Simulations (Rank 1 - PC1)}

\bigskip

Dear Members of the Selection Committee,

As a State Engineer in Artificial Intelligence from ESI-SBA, I am writing to express my enthusiastic application for the PhD position in "Interpretable and Robust Machine Learning for Physics-Based Numerical Simulations" within the Engineering Digital Twin (EDT) national program. This application is motivated by the close alignment between my technical background in hybrid physics-data modeling and the research objectives of the Catalyst project.

Currently, I am completing my thesis at Sonatrach, where I have developed a **Physics-Informed Digital Twin** for gas turbine fleet monitoring. My research focuses on using machine learning to emulate and accelerate numerical solvers for complex physical problems—a direct parallel to the objectives of the **Catalyst (Reliable Hybrid Model Forge)** project. Specifically, I have designed custom **PINN (Physics-Informed Neural Network)** loss functions that enforce thermodynamic consistency within generative surrogate models based on **Mamba-SSM** and **xLSTM** architectures.

The core challenge of your project—addressing the "black-box" nature of ML surrogates and their sensitivity to out-of-distribution data—is exactly what motivates my current work. In my Sonatrach project, I have encountered the critical need for **quantitative indicators of model reliability** when deploying surrogates on edge devices. I am particularly eager to apply advanced statistical methods, such as **Sensitivity Analysis**, **Optimal Transport**, and **Statistical Depth**, to define the reliability domains of these hybrid models. 

My background in high-dimensional statistical analysis and my experience with production-grade PyTorch and C++ implementations provide me with the necessary tools to contribute to the **Artemis platform** and the broader EDT ecosystem. I am a highly autonomous researcher, comfortable bridging the gap between theoretical statistics and industrial physical simulations.

Thank you for your time and consideration. I look forward to the possibility of discussing how my experience in physics-aware AI and industrial surrogates can contribute to the success of the Catalyst project.

Sincerely,

\bigskip

\textbf{Ismail AISSAOUI}

\end{document}
